\section{国际法上的空间}
\begin{enumerate}
  \item 国际法内容最多的一章。
  \item 领土制度。构成。陆水空底土。
        \begin{enumerate}
          \item 陆 \yjt 边境制度。便利；相邻权；界标\zd{双方代表在场}。
          \item 水 \yjt 河流制度。河流、湖泊、3 大海域：内海、领海、群岛水域。
                \begin{enumerate}
                  \item 界河。主航道 \hz 中心线为界；可以航行但不可停靠；修建设施须经对方许可。
                  \item 多国河流。仅对沿岸国开放；主权归属流经国；不得改道 \hz 堵塞河流。
                  \item 国际河流。非军用船舶无害通过；主权归属流经国；条约成立专门机构管理。
                  \item 国际运河。所在国完全主权；条约确立航行制度；一般对所有国家开放。
                \end{enumerate}
          \item 空 \yjt 完全主权，\zd{劫机犯或引渡或起诉}。
          \item 底土 \yjt 完全主权。自然资源。
        \end{enumerate}

        \begin{kk}{取得方式}
          \begin{enumerate}
            \item 先占（一定是国家的行为）。天然形成的无主地。\zd{\bdy 最先发现，最先开发}。解决历史问题。
            \item \zd{时效}。有争议，我国不承认其合法性。
            \item 添附。自然、人工。\zd{以领土为基础}，不能侵害相邻权。
            \item 征服。\zd{武力政府违法}。
            \item 割让（看个案）。条约是否自愿平等。\zd{强制则不合法}。
          \end{enumerate}
          \tcblower
          新的两种。殖民地独立；公民投票看国内法有无协议，相关国能否达成协议。个案分析。
        \end{kk}

        \begin{kk}{限制}
          \begin{enumerate}
            \item 一般限制。领海无害通过；外交特权与豁免……
            \item 特殊。\zd{共管、租界、国际地役、势力范围。效力取决于据以产生的条约是否合法}。
          \end{enumerate}
        \end{kk}

  \item 海洋制度。一定会考。海洋水域。海洋底土。放图。
        \begin{enumerate}
          \item 内\zd{水} \yjt 内\zd{海} \yjt \zd{领}海 \yjt \zd{毗}连区 \yjt \zd{专属}经济区 \yjt \zd{大}陆架 \yjt \zd{公}海/国际海域。
          \item 内海。\zd{完全主权}，非请勿入。\zd{渤海、琼州海峡}属于内海。
                \begin{enumerate}
                  \item 港口刑事管辖三主动。\zd{扰乱港口安宁；案情重大；本国人受害}。船旗国领事 \hz 船长提被动管。
                \end{enumerate}
          \item 领海。无害通过权。不允许军舰；连续不停迅速通过；潜水器浮出水面挂出国旗。
          \item 毗连区。\zd{海关（走私）、财政（逃税）、移民（偷渡）、卫生（检验检疫）、国家安全（我国）}。
          \item 专属经济区。自然资源的专属勘探开发和管辖；事先报备；宣告有效。
          \item 公海。
                \begin{enumerate}
                  \item 船旗国。一船一旗，中途变更、悬挂多国国旗视为无国旗。
                  \item 普遍管辖。\zd{海盗、非法广播、贩毒、奴隶}。
                  \item 登临权/临检权（起点公海）vs 紧追权（其他延伸到公海）。对象只能非政府、国家船舶。
                  \item 紧追前须发出停止信号（\zd{敲锣打鼓}）、不能中断（交接班算中断）、\zd{追到他国领海}。
                \end{enumerate}
          \item 海洋底土。内海、领海、专属经济区下覆底土。大陆架；国际海底区域平行开发（勘探者和国际海底管理局）。
        \end{enumerate}

        \begin{kk}{大陆架}
          \begin{enumerate}
            \item 宽度 \zd{200 到 350}。举证获得（不以占领 \hz 公告为依据）。
            \item 自然资源（非生物资源）专属开发和管辖权。
            \item 自愿开发，内外有别。200 海里以外\zd{国际海底管理局}缴费，\zd{发展中国家}某些条件下享有免缴权。
          \end{enumerate}
        \end{kk}

        \begin{kk}{群岛水域}
          \begin{enumerate}
            \item 岛屿封闭线以外，群岛基线以内，内水和领海之间，主轮廓，不能隔断其他国家领海 \hz 专属经济区。
            \item 领水 = 内水 + 领海，完全主权，
            \item 无害通过制度。船无害通过；飞机经许可走空中通道（是权利）。
          \end{enumerate}
        \end{kk}

        \begin{kk}{国际海峡}
          \begin{enumerate}
            \item 两端连接的是公海 \hz 专属经济区。
            \item 通行制度。最最重要。
                  \begin{enumerate}
                    \item \zd{过境通行制}。飞机、船舶不经许可、不打招呼、无害通过。
                    \item （船）\zd{无害通过权}。一国大陆和该国岛屿构成的国际海峡且该岛屿\zd{向海}一面\zd{有}方便的\zd{航道}。
                    \item 如果够宽中间航道，\zd{自由通行制}，可走走停停。
                    \item 沿岸国可以\zd{指定航道、规定分道分航}（是权利），但须征得国际海事组织同意，危险情况公布。
                  \end{enumerate}
          \end{enumerate}
        \end{kk}
  \item 特殊空间法律制度。
        \begin{kk}{南极}
          可考性相对差一些。
          \begin{enumerate}
            \item 人类共同利益原则、和平目的原则、自由科考原则。
            \item \zd{冻结}南极领土主权要求。不前进不后退。
          \end{enumerate}
        \end{kk}

        \begin{kk}{外空}
          航空器能到达的高度。
          \begin{enumerate}
            \item 登记制度。联合国秘书长处登记。法律后果：是管辖权和所有权的依据。
            \item 营救制度。\zd{援助、通知、送回}。
            \item 责任制度。\zd{发射国扩大解释}。
                  \begin{enumerate}
                    \item 空对地绝对责任，空对空过错责任。
                    \item 单方过错单方承担，都有过错按比例，没有过错损失自担。
                    \item 不调整\zd{缔约国和本国国民，缔约国和受邀参与发射活动的法国人}。
                  \end{enumerate}
          \end{enumerate}
        \end{kk}
  \item 国际环保法。
        \begin{enumerate}
          \item 控制危险废物越境转移/巴塞尔公约。
                \begin{enumerate}
                  \item \zd{缔约国}之间。
                  \item 特定进口\zd{书面}同意。
                  \item 有无害环境的处置方法。\al{核废水}。
                  \item 必须要\zd{有保险 \hz 担保}。
                \end{enumerate}
          \item 大气环保。92 年联合国气候变化框架公约、97 年京都议定书、15 年巴黎协定。重点热点。
                \begin{enumerate}
                  \item \zd{共同但有区别责任原则}。
                  \item 工业化国家、发达国家（\zd{拿出资金和技术援助}）、发展中国家。
                  \item 净排量计算方式。实际排量扣掉本土森林所能吸收的二氧化碳数量。
                  \item 3 种折抵方式。\zd{向发达国家购买额度；向发展中国家输出绿色技术；集团方式（欧盟）}。
                  \item 2020 年以后，国家自主贡献方式，有减排义务，但不再适用公约的强制摊派指标。
                \end{enumerate}
          \item \zd{濒危野生动植物种国际贸易公约}。\zd{没考过，大纲新增}。以清单为基础的许可证制度。
                \begin{enumerate}
                  \item 活体和标本。
                  \item 附件一。要求取得进出口许可证。
                  \item 附加二。要求取得出口许可证。
                  \item 附件三。缔约国自行决定、管理。
                \end{enumerate}
        \end{enumerate}
\end{enumerate}
